\documentclass[conference]{IEEEtran}
\IEEEoverridecommandlockouts
% The preceding line is only needed to identify funding in the first footnote. If that is unneeded, please comment it out.
%Template version as of 6/27/2024

\usepackage{cite}
\usepackage{amsmath,amssymb,amsfonts}
\usepackage{algorithmic}
\usepackage{graphicx}
\usepackage{textcomp}
\usepackage{xcolor}
\def\BibTeX{{\rm B\kern-.05em{\sc i\kern-.025em b}\kern-.08em
T\kern-.1667em\lower.7ex\hbox{E}\kern-.125emX}}
\begin{document}

\title{Roole at the SMT-COMP 2026}

\author{\IEEEauthorblockN{Jan Onderka}
  \IEEEauthorblockA{%
    \textit{University of Freiburg}\\
    Freiburg, Germany \\
  onderka@cs.uni-freiburg.de}
  \and
  \IEEEauthorblockN{Armin Biere}
  \IEEEauthorblockA{%
    \textit{University of Freiburg}\\
    Freiburg, Germany \\
  biere@cs.uni-freiburg.de}
}

\maketitle

\begin{abstract}
  We present Roole, a Satisfiability Modulo Theories (SMT)
  solver for the theory of quantifier-free bitvectors,
  as submitted to the SMT-COMP 2026.
\end{abstract}

\section{Introduction}

\textbf{Roole} \cite{roole} is a new SMT solver for the theory of quantifier-free bitvectors
based on simplified Input-based Three-Valued Abstraction Refinement \cite{onderka2026}
combined with DPLL-style search with learning and backjumping.
\textbf{Roole} computes bitvector operations in three-valued bitvector abstraction \cite{onderka2022},
using no bitblasting nor any underlying SAT solvers.
It is dual-licensed under both the MIT and Apache-2.0 licence.

\textbf{Roole} is written in Rust and is designed with focus on certified solving,
with the proof-checker \textbf{Roolean} \cite{roolean}
written in Lean that can check that the results are sound by essentially solving once again
with the aid of \emph{oracle certificates} from \textbf{Roole} that guide its choices for better performance,
the soundness of \textbf{Roolean} results proven through Lean theorems.
However, as SMT-COMP does not require results to be certified, we submit \textbf{Roole} as a standalone solver.

\section{Approach}
We based \textbf{Roole} on our previous work on Input-based Three-Valued Abstraction Refinement \cite{onderka2022},
using the fact that SMT solving of quantifier-free formulas is similar to model-checking the property
\textbf{EX}[$p$] (there exists a path where $p$ holds in the next state), simplifying to the following.

Consider a function $f: V_1 \rightarrow V_2$. Introduce abstract values in sets $\hat{V}_1$ and $\hat{V}_2$
related by the \textit{concretization function} $\gamma: \hat{V}_\Box \rightarrow 2^{V_\Box}$
that determines the concrete values overapproximated by the abstract value.
The function $\hat{f}: \hat{V}_1 \rightarrow \hat{V}_2$ soundly overapproximates $f$ exactly if
\begin{equation}
  \forall \hat{v} \in \hat{V}_1 \; . \; \forall v \in \gamma(\hat{v}) \; . \;
  f(v) \in \gamma(\hat{f}(\hat{v})).
\end{equation}
Once sound abstract functions are implemented for basic operations,
they can be composed while retaining soundness. Finally, considering $f$ as the
term determining satisfaction of the formula by a given assignment $v \in V_1$,
the result of checking satisfiability is equal to
\begin{equation}
  \exists v \in V_1 \; . \; f(v).
\end{equation}
Brute-force computation using $f$ can then be replaced by computation using its sound abstraction $\hat{f}$
on a set of abstract assignments $\hat{A} \subseteq \hat{V}_1$
provided that no spurious abstract values are considered
($\forall \hat{a} \in \hat{A} \; . \; \gamma(\hat{a}) \neq \emptyset$)
and all concrete values are considered
($\left(\bigcup_{\hat{a} \in \hat{A}} \gamma(\hat{a})\right) = V_1$).

\section{Implementation}

The efficiacy of the approach depends on balancing the computation speed of abstract operations
and the number of abstract assignments needed to consider.

Currently, we use the three-valued bitvector domain
with fast computation of arithmetic operations \cite{onderka2022}.
To search through the space of the abstract assignments for those where satisfiability
or unsatisfiability of all overapproximated concrete assignments is shown,
we use a search algorithm in the style of a Davis-Putnam-Logemann-Loveland (DPLL)
procedure with additional basic machinery for backjumping and learning of abstract assignments
found to be unsatisfiable, inspired by Conflict-Driven Clause Learning (CDCL).

We note that the current implementation is a first stepping stone, without the use of commonly used
techniques such as preprocessing or sophisticated CDCL techniques, more powerful domains,
or any fine-tuning.

\bibliographystyle{IEEEtran}
\bibliography{main}

\end{document}
